\scanpage{14}
\mspage{10}
% source: top, complementary polygon
Ceci posé, on définit le \uline{polygone complémentaire} ${\unsure{\Phi|\varphi}}^{*}$ de $\Phi$ comme
\[
\begin{aligned}
\Phi^*&=\widehat\Phi\setminus\Phi=\text{ens des }\deleted{\text{intervalles}}\inserted{\text{strates}}\\
&\qquad\text{interstitielles généralisées de }\delta\Phi,\text{ qui ne}\\
&\qquad\text{sont pas }\in\Phi.
\end{aligned}
\]

% source: middle, redundant vertices
On aura
\[
\delta\Phi^*=\delta\Phi\setminus\Phi_{\deleted{\unsure{\mathrm{le}}}\mathrm{red}},\quad\text{soit }T'
\]
où $\Phi_{\mathrm{réd}}$ est l'ens. des sommets redondants de $\Phi$, i.e.
\[
\Phi_{\mathrm{red}}=\{s\in\Phi_{\unsure{0|o}}\mid\operatorname{ord}(s\unsure{.|,}\Phi)=2\}
\]

% source: middle-lower, digression
\deleted{On appelle}

Il y a lieu ici de faire la petite digression sur la nature des sommets d'un $\unsure{\widehat M|\widehat{\mathscr M}}$-polygone (cinq types: sommets isolés, sommets-bord propres et impropres, sommets internes propres (ou redondants) ou impropres (dits lacunaires)) -- cf. pages \unsure{110|10}, \unsure{111|11}. On trouve \deleted{\unsure{aussitôt}} que $\Phi^*$ n'a pas de sommets redondants, donc $\delta\Phi^{\unsure{\circ *|**}}=\delta\Phi^*$.

% source: bottom, double complement
On obtient
\[
\begin{aligned}
\deleted{\unsure{\delta}}\Phi^{**}&=\text{ens des strates interstitielles}\inserted{X}\text{ de }T'=\delta\Phi\setminus\Phi_{\mathrm{red}}\\
&\quad\deleted{\text{\unsure{pour lesquels il existe}}}\inserted{\text{telles qu'il existe}}Y\in\Phi\ \unsure{\text{avec}|\text{tel que}}\ Y\ll X.
\end{aligned}
\]
On \unsure{sait|voit} que $\Phi^{\unsure{**}}$ s'obtient à partir de $\Phi$
% archive sheet number: 13
